\documentclass{article}
\usepackage{graphicx}
\usepackage{amsfonts,amssymb,amsmath} 
\usepackage{lmodern,adjustbox}
\usepackage{float}
\usepackage{physics}
\usepackage{tikz}
\usetikzlibrary{quantikz}
\usepackage{mathtools}
% \usepackage[a4paper, left = 3.5cm, right = 3.5cm, top = 3.5cm, bottom = 3.5cm ]{geometry}
% \usepackage[a4paper, left = 2cm, right = 2cm, top = 2cm, bottom = 2cm ]{geometry}
\usepackage[a4paper, left = 1cm, right = 1cm, top = 2cm, bottom = 2cm ]{geometry}
\usepackage{listings}

\newtheorem{definition}{Definition}[section]
\newtheorem{proof}{Proof}[section]

\begin{document}
	
	\begin{titlepage}
		
		\title{Report of Introduction to Quantum Computing \\ Quantum Error Correction}
		\author{Student: Diego Arcelli\\ Student ID: 647979 \\
			Professors: Gianna Maria Del Corso,\\Anna Bernasconi, Roberto Bruni}
		\date{Academic Year 2021-2022}
		\maketitle
		\centering
		\includegraphics[width=10cm]{./images/unipi_logo.png}
		
	\end{titlepage}
	\tableofcontents
	\newpage
	
	\section{Quantum Error Correction}
	 
	The hardware devices used to implement quantum circuits are sensitive to noise (much more than classical computers). This might cause the information carried by the circuit to be modified in an unpredictable way. For this reason being able to detect and correct these errors it's a topic of major importance. \\
	Sadly handling errors in the quantum case it's more troublesome than classical computing: first, we cannot copy arbitrary qubits (because of the no cloning theorem), then we cannot measure a qubit to check its state because this would destroy the superposition and, moreover, errors are continuous (while in classical computing we only have bits flip errors).
	
	\section{General Framework}
	We start by describing a general framework for quantum error correction, and in the next section we'll se how it can be applied in practice. The first step to protect a quantum state from a possible error is to define an error model, which is a description of the type of error it might happen. Once we have the error model, in order to defend a quantum state $\ket{\psi}$ from the error that might occur, we'll define an encoding operator $U_{enc}$ which, by adding some extra ancillary qubits to the state, will produce an encoded state $\ket{\psi_{enc}}$:
	\[ U_{enc}\ket{\psi}\ket{00\dots0} = \ket{\psi_{enc}} \]
	The space of all the possible encoded states will form a sub-space of the Hilbert space that we'll call the Quantum Error Correcting Code (QECC). If $\ket{\psi}$ is composed by $k$ qubit, and with the encoding we add $n$ qubits, then the code $\mathcal{C}$ will be a sub-space of the $2^{k+n}$ Hilbert space. The idea of encoding the state in a higher dimensional state, is to use the additional qubits added with the encoding to store information about the error that occurred on the state. Then we define a recovery operator $\mathcal{R}$ that, if an error occurs on the encoded state $\ket{\psi_{enc}}$, will undo the effect of the error, based on the information gained with the error syndrome. Finally we'll decode $\ket{\psi_{enc}}$ back to $\ket{\psi}$ with a decoding operator:
	\[U_{dec}\ket{\psi_{enc}} = \ket{\psi}\ket{0\dots0}\]
	In general given a quantum code $\mathcal{C} \subset \mathcal{H}$, we say that a noise operator $\mathcal{E}$ is recoverable it there exists a recovery operator $\mathcal{R}$ such that for every density matrix $\rho$ associated with a state $\ket{\psi} \in \mathcal{C}$:
	\[ \mathcal{R(\mathcal{E}(\rho))} = \rho\]
		
	\section{Correcting single qubit errors}
	In the following sections we'll analyze some of the most relevant type of errors that can occur on a single qubit, and we'll show circuits which can detect and correct the errors, using the framework explained before.	
	\subsection{Correcting bit flip errors}
	The simplest error that we might have is the bit flip error, where a generic qubit in state $\alpha\ket{0} + \beta\ket{1}$ it's transformed in $\alpha\ket{1} + \beta\ket{0}$. This can be seen as an application of the $X$ Pauli's gate to the qubit. The idea of how to correct this error comes form classical computing, where if we have a noisy channel where a classical bit can flip with probability $p$, to reduce the probability of receiving the wrong bit, we can repeat the bit three times. In this way the receiver can decode the string of three bits simply by choosing the majority bit, so, for instance, if it receives 101, it will decode it as 1. Note that if two or all three bits get flipped, the receiver will decode the wrong bit, but this happens with probability $3p^2 + p^3$ instead of $p$. We can apply the same idea in the quantum case, by choosing the following encoding for the states $\ket{0}$ and $\ket{1}$:
	\[ \alpha\ket{0} + \beta\ket{1} \rightarrow \alpha\ket{000} + \beta\ket{111}\]
	We say that the logical states $\ket{0}$ and $\ket{1}$ have been encoded in the physical states $\ket{000}$ and $\ket{111}$. We want to apply the same strategy we used for the classical case, so checking if the values of the three qubits are the same, and if not, we impose the value of the majority qubit. Of course we cannot do that by measuring the state, otherwise the superposition will be destroyed. What we can do instead is using some extra qubits, to store in them the information of which qubit of the state has been flipped, and we'll call these extra qubits syndrome qubits. The full circuit to get these operations done is the following:
	\begin{figure}[H]
	\centering
	\resizebox{10cm}{!}{%
	\begin{quantikz}
		\lstick{$\ket{\psi}$} & \ctrl{1}\gategroup[3,steps=2,style={dashed,
			rounded corners,fill=blue!10, inner xsep=2pt},
		background]{{\sc Encoding}} & \ctrl{2} & \qw & \ctrl{3}\gategroup[5,steps=4,style={dashed,
			rounded corners,fill=blue!10, inner xsep=2pt},
		background]{{\sc Detection}} & \qw & \qw & \qw & \qw & \gate{X^{\lnot x \land y}}\gategroup[3,steps=1,style={dashed,
			rounded corners,fill=blue!10, inner xsep=2pt},
		background]{{\sc Correction}} & \qw & \ctrl{2}\gategroup[3,steps=2,style={dashed,
			rounded corners,fill=blue!10, inner xsep=2pt},
		background]{{\sc Decoding}} & \ctrl{1} & \qw\\
		\lstick{$\ket{0}$}  & \targ{} & \qw & \qw & \qw & \ctrl{2} & \ctrl{3} & \qw & \qw & \gate{X^{x \land y}} & \qw & \qw & \targ{} & \qw\\
		\lstick{$\ket{0}$}  & \qw & \targ{} & \qw & \qw & \qw & \qw & \ctrl{2} & \qw & \gate{X^{x \land \lnot y}} & \qw & \targ{} & \qw & \qw\\
		\lstick{$\ket{0}$} & \qw & \qw & \qw & \targ{} & \targ{} & \qw & \qw & \meter{}\gategroup[2,steps=2,style={dashed,
			rounded corners,fill=blue!10, inner xsep=2pt},
		background, label style={label position=below,anchor=
			north,yshift=-0.2cm}]{{\sc Syndrome}} & \cw \ x &  &  &  \\
		\lstick{$\ket{0}$} & \qw & \qw & \qw & \qw & \qw & \targ{} & \targ{} & \meter{} & \cw \ y &  &  & 
	\end{quantikz}
	}
	\caption{Bit flip error correction circuit}
	\end{figure}
	\noindent The first three qubits are used for the encoding, while the other two are the syndrome qubits. The first two CNOT gates are used to apply the encoding. Then, with the second two CNOT gates, we check if the first two qubits are equal. If they are equal and they're both $\ket{0}$ then the first syndrome qubit is not flipped. If they are equal an they're both $\ket{1}$ the first syndrome qubit it's flipped twice (so it goes back to 0). If they're different the syndrome qubit is flipped once and so it is set to $\ket{1}$. With the last two CNOT gates we repeat the same procedure, but this thy we compare the second and the third qubits and we act on the second syndrome qubit. Finally we measure the syndrome qubits and we get the values $x$ and $y$. The value of $x$ tells us if the first and the second qubit were the same, while the value of $y$ tells us if the second and the third qubit were the same. Like in the classical case, if a bit flip error occurs in 2 or all the 3 qubits, then the circuit will decode the wrong state $\alpha\ket{1} + \beta\ket{0}$, but anyway the circuit reduces this probability. 

	\subsection{Correcting phase flip errors}
	Another type of error which does not have an equivalent in the classical case is the phase flip error, where a generic qubit $\alpha \ket{0} + \beta\ket{1}$ is transformed in $\alpha\ket{0} - \beta\ket{1}$. This can be seen as an application of the $Z$ Pauli gate. It turns out that solving phase flip errors is very similar to solve bit flip errors, since phase flip errors in the canonical basis become bit flip errors in the Hadamard basis since $X = HZH$ and because of that, $Z\ket{+} = \ket{-}$ and $Z\ket{-} = \ket{+}$. We can exploit this fact by using the following encoding:
	\[ \alpha\ket{0} + \beta\ket{1} \rightarrow \alpha\ket{+++} + \beta\ket{---}\]
	which can be obtain from the bit flip encoding, by applying to each one of the three qubit the $H$ gate. Suppose that with this encoding a phase flip error occurs in the second qubit. Then, since in the Hadamard basis, a phase flip acts like a bit flip, the encoded state becomes:
	\[  \alpha\ket{+-+} + \beta\ket{-+-} \]
	if we now apply and $H$ gate to each qubit of the state, it becomes:
	\[  \alpha\ket{010} + \beta\ket{101} \]
	which is like if a bit flip error occurred in the second bit of the encoding for the bit flip error. Therefore we can simply run the circuit, for the bit flip correction, and we'll correctly recover the initial phase flip error. So, like the bit flip circuit, the circuit is able to correct a phase flip error only if it occurs in jut one of the three qubits. 
	

	\subsection{Nine-qubits correcting code}
	We have seen a circuit which can correct bit flip errors and a circuit which can correct phase flip errors, but both of them fail if a phase flip and a bit flip error happens on the encoded state at the same time. The intuition is to start from the encoding which protects from phase flip errors:
	\[ \alpha\ket{0} + \beta\ket{1} \rightarrow \alpha\ket{+++} + \beta\ket{---}\]
	and apply to each one of the three qubit the circuit for the bit flip encoding, in order to protect each qubit from bit flip errors, ending up with the following nine-qubit encoding:
	\[ \alpha\ket{0} + \beta\ket{1} \rightarrow \frac{\alpha}{\sqrt{8}}(\ket{000} + \ket{111})^{\otimes 3}+ \frac{\beta}{\sqrt{8}} (\ket{000} - \ket{111})^{\otimes 3}\]
	So we have three block of three qubits each. Within each block, all the qubit are supposed to be the same, so we can run the bit flip error correction circuit in every one of the block to protect each block by bit flip errors. Now, if we consider the thee block, the sign of every block should be the same, but if there's been a phase flip error in one of the block, the sign on that block will be flipped. The idea is the same but this time we cannot apply the circuit we saw before, but what we can do is use this circuit:
	\begin{figure}[H]
	\centering
	\resizebox{8cm}{!}{%
	\begin{quantikz}
		& \gate{H} & \ctrl{9} & \qw & \qw & \qw & \qw & \qw & \qw & \qw & \qw & \qw & \qw & \qw & \gate{H} & \qw & \gate{Z^{\lnot x \land y}} & \qw \\
		& \gate{H} & \qw & \ctrl{8}  & \qw & \qw & \qw & \qw & \qw & \qw & \qw & \qw & \qw & \qw & \gate{H} & \qw & \qw & \qw\\
		& \gate{H} & \qw & \qw & \ctrl{7} & \qw & \qw & \qw & \qw & \qw & \qw & \qw & \qw & \qw & \gate{H} & \qw & \qw & \qw \\
		& \gate{H} & \qw & \qw & \qw & \ctrl{6} & \qw & \qw & \ctrl{7} & \qw & \qw & \qw & \qw & \qw & \gate{H} & \qw & \gate{X^{x \land y}} & \qw\\
		& \gate{H} & \qw & \qw & \qw & \qw & \ctrl{5} & \qw & \qw & \ctrl{6} & \qw & \qw & \qw & \qw & \gate{H} & \qw & \qw & \qw \\
		& \gate{H} & \qw & \qw & \qw & \qw & \qw & \ctrl{4} & \qw & \qw & \ctrl{5} & \qw & \qw & \qw & \gate{H} & \qw & \qw & \qw\\
		& \gate{H} & \qw & \qw & \qw & \qw & \qw & \qw & \qw  & \qw & \qw & \ctrl{4} & \qw & \qw & \gate{H} & \qw & \gate{X^{x \land \lnot y}} & \qw\\
		& \gate{H} & \qw & \qw & \qw & \qw & \qw & \qw & \qw & \qw & \qw & \qw & \ctrl{3} & \qw & \gate{H} & \qw & \qw & \qw\\ 
		& \gate{H} & \qw& \qw & \qw & \qw & \qw & \qw & \qw & \qw & \qw & \qw & \qw & \ctrl{2} & \gate{H} & \qw & \qw  & \qw\\
		& \lstick{$\ket{0}$} & \targ{} & \targ{} & \targ{} & \targ{} & \targ{} & \targ{} & \qw & \qw & \qw & \qw & \qw  & \qw & \meter{} & \cw \ x \\
		& \lstick{$\ket{0}$} & \qw & \qw & \qw & \qw & \qw & \qw & \targ{} & \targ{} & \targ{} & \targ{} & \targ{} & \targ{} & \meter{} & \cw \ z
	\end{quantikz}
	}
	\end{figure}  \[ \ket{0} \rightarrow \ket{00000} + \ket{00000} + \ket{00000} + \ket{00000} + \ket{00000} + \ket{00000} + \ket{00000} + \ket{00000} + \]
	\[ \ket{00000} + \ket{00000} + \ket{00000} + \ket{00000} + \ket{00000} + \ket{00000} + \ket{00000} + \ket{00000} \]
	\noindent it rather messy, but what the circuit does is comparing the sign of the first block with the sign of the second block and if they're different set the first ancillary qubit to 1, and then it does the same thing with the second and third block, so that we'll measure a syndrome for each possible error.\\
	
	To show why the circuit works, after the application of $H^{\otimes 9}$ the blocks will be transformed like that:
	\[  (\ket{000} + \ket{111})^{\otimes 3} \rightarrow 2(\ket{000} + \ket{011} + \ket{101} + \ket{110})^{\otimes 3} = \ket{\psi^+}^{\oplus 3}\]
	\[  (\ket{000} - \ket{111})^{\otimes 3} \rightarrow 2(\ket{001} + \ket{010} + \ket{100} + \ket{111})^{\otimes 3} = \ket{\psi^-}^{\oplus 3}\]
	\[\alpha \ket{0} + \beta\ket{1} \rightarrow \alpha\ket{\psi^+}\ket{\psi^+}\ket{\psi^+} + \beta\ket{\psi^-}\ket{\psi^-}\ket{\psi^-}\]
	showing all the steps would be too messy and not even interesting, but the idea is that now each qubit of each block is composed a series of binary string where for the encoding of $\ket{0}$ each string is composed by an even number of 1s, while for the encoding of $\ket{0}$ by an odd number of 1s. The circuit exploit this property, since if there's been no phase error, the first three CNOT gates will flip the first ancillary qubit an even number of times for the encoding of $\ket{0}$, and an odd number of times for the encoding of $\ket{1}$. The the second three CNOTs will flip the first qubits again an even number of times for the encoding of $\ket{0}$, and an odd number of times for the encoding of $\ket{1}$. Since an even number plus an even number is an even number, and an odd number plus an odd number is an odd number, for both the encoding of $\ket{0}$ and $\ket{1}$ it gets flipped an even number of times and so it remains $\ket{0}$.\\
	Suppose that there's been a phase flip error on the first block. Then state after the application of the Hadamard gates will be:
	\[ \alpha\ket{\psi^-}\ket{\psi^+}\ket{\psi^+} + \beta\ket{\psi^+}\ket{\psi^-}\ket{\psi^-}\]
	so this time when we apply the first three CNOTs for the encoding of $\ket{0}$ the first ancillary qubit will be flipped an odd number of times, and an even number of times for the encoding of $\ket{1}$. Then with the second three CNOTs the ancillary qubits will be flipped an even number of times for the encoding of $\ket{0}$ and an odd number of times for the encoding of $\ket{1}$. Since an even number plus and odd number is an odd number, the first ancillary is flipped in total and odd number of times for both the encoding of $\ket{0}$ and $\ket{1}$. The same reasoning can be applied with the second and third triples of CNOTs and with the third and the fourth.
	
	\subsection{Correcting generic errors}
	Since now we only saw bit flip and phase flip errors, but how do we deal with other type of errors, like a generic rotation of the qubit in the Bloch sphere. Consider a generic error on a single qubit defined by a matrix $E \in \mathbb{C}^{2 \times 2}$, and consider the four Pauli matrices $I$, $X$, $Y$ and $Z$. As we already saw $X$ and $Z$ represent respectively bit flip and phase flip errors. Since $Y = iXZ$, it can be seen as bit flip error followed by a phase flip error (the global phase factor $i$ does not matter), while the gate $I$ can be simply seen as the no-error gate. These four matrices span the space of the $2 \times 2$ complex matrices, so the before error matrix $E$ can be written as:
	\[ E = \alpha_0I + \alpha_1X + \alpha_2Y + \alpha_3Z, \quad \alpha_0, \alpha_1, \alpha_2, \alpha_3 \in \mathbb{C}\]
	Basically we've rewritten $E$ as a linear combination of bit flip and phase flip errors, which we already know how to solve. This is important because there's a theorem which state that if a code $\mathcal{C}$ corrects a set of errors $\mathcal{E} = \{E_i\}$, then $\mathcal{C}$ corrects also the span of $\mathcal{E}$. So if we're able to correct bit flip and phase flip errors, then we're able to correct any kind of error.
	
	\section{Stabilizer codes}
	It turns out that the codes that we saw so far and many others can be described using a particular formalism, introduced by Gottersman in \cite{Gottesman1997StabilizerCA}, which makes it easier to define quantum error correcting codes. In the following section we'll introduce the stabilizer formalism and we'll see how we can apply it to quantum error correction.
	
	\subsection{Stabilizer formalism}
	We say that a unitary operator $U$ stabilizes the state $\ket{\psi}$ if $U\ket{\psi} = \ket{\psi}$, which means that $\ket{\psi}$ is an eigenstate of $U$ with eigenvalue $1$. If $V$ and $U$ stabilize $\ket{\psi}$ the also their product and their inverses stabilize $\ket{\psi}$. Since matrix multiplication is an associative operation, this means that the set of unitary matrices that stabilizes $\ket{\psi}$ form a group. 	
	Consider the set of Pauli operators $\Pi = \{I,X,Y,Z\}$, we then define $\Pi^n$ as:
	\[\Pi^n = \{A_1\otimes \dots \otimes A_n\ \colon A_j \in \Pi\ \forall j\}\]
	which is the set $2^n\times2^n$ matrices which can be obtained by the tensor product of $n$ Pauli matrices. A stabilizer $\mathcal{S}$ is an abelian subgroup of $\Pi^n$. \\
 	Given a stabilizer $\mathcal{S}$ we can define the code-space generated by $\mathcal{S}$ as $\mathcal{T}(\mathcal{S}) = \{ \ket{\psi} \in \mathcal{H} \colon M\ket{\psi} = \ket{\psi }\ \forall M \in \mathcal{S}\}$ which is the set of state stabilized by all the elements of $\mathcal{S}$. We can interpret $\mathcal{T}(\mathcal{S})$ as a QECC where the states in $\mathcal{T}(\mathcal{S})$ are the encoded states.
	
	\subsection{Thee and nine qubits encoding revisited}
	Let's now revisit the three qubit code for bit flip correction. What we were trying to do was check if the first and the second qubits where the same, and the same thing for the second and the third qubits. We can equivalently say that we where verifying if the encodes state has eigenvalue $1$ with respect to the operators $Z_1 = Z\otimes Z \otimes I$ and $Z_2 = I \otimes Z \otimes Z$, since:
	\[(Z\otimes Z \otimes I)(\alpha\ket{000} + \beta\ket{111}) = \alpha\ket{000} + \beta\ket{111}\]
	if instead there's a bit flip error in the first qubit we get that:
	\[(Z\otimes Z \otimes I)(\alpha\ket{010} + \beta\ket{101}) = -(\alpha\ket{010} + \beta\ket{101}) \]
	so the eigenvalue of the state wit respect to that operator becomes $-1$. So what we were doing was checking if the encoded state has eigenvalue 1 with respect both of $Z_1$ and $Z_2$, while if for $Z_1$ has eigenvalue $-1$ it means that the first and the second qubits are different, while if it has eigenvalue $1$ with $Z_2$, it means that the second and the second qubits are different. It's easy to check that we could have done the same bit flip correcting procedure, by comparing first the first and the second qubits, and then the first and the third, the syndrome would have been different but it could be used to correct the error. In therm of stabilizers, this corresponds to using the matrices $Z_1$ (the same we used before) and $Z_3 = Z \otimes I \otimes Z$, since also for $Z_3$ we have $Z_3 \ket{\psi_{enc}} = \ket{\psi_{enc}}$. In fact the set of stabilizers of $\ket{\psi_{enc}}$ is $\mathcal{S} =\{Z_1, Z_2, Z_3\}$, but we can only use $Z_1$ and $Z_2$ because they are generators of $Z_3$, since $Z_1Z_2 = Z_2Z_1 = Z_3$. \\
	Another way to put this is by saying that the correcting code must have eigenvalue 1 with respect of the operators:
	\[
	\begin{array}{c|c c c}
		g_1 & Z & Z & I \\
		g_2 & I & Z & Z \\
	\end{array}
	\]
	The same reasoning applies to the phase flip correction but with the matrices $X_1 = X\otimes X \otimes I$ and $X_2 = I \otimes X \otimes X$, and we won't show it since it's straightforward to verify that is true. Instead it is more interesting to see. In instead we think about the nine qubit encoding the corresponding stabilizer representation is the following:
	\[
	\begin{array}{c|c c c c c c c c c }
		g_1 & Z & Z & I & I & I & I & I & I & I \\
		g_2 & I & Z & Z & I & I & I & I & I & I \\
		g_3 & I & I & I & Z & Z & I & I & I & I \\
		g_4 & I & I & I & I & Z & Z & I & I & I \\
		g_5 & I & I & I & I & I & I & Z & Z & I \\
		g_6 & I & I & I & I & I & I & I & Z & Z \\
		g_7 & X & X & X & X & X & X & I & I & I \\
		g_8 & I & I & I & X & X & X & X & X & X \\
	\end{array}
	\]
	the first size basically correspond to the bit flip stabilizer applied to each group of three qubits. 
	
	\subsection{Errors with stabilizer codes}
	 Consider an error matrix $E \in \Pi^n$. $E$ can either commute or anti-commute with the operators in $\mathcal{S}$. If $\forall M \in \mathcal{S}\ ME = -EM$ then $ME\ket{\psi} = -EM\ket{\psi} = -E\ket{\psi}$, thus the error defined by $E$ can be detected. If instead $\forall M \in \mathcal{S}\ ME = EM$ then $ME\ket{\psi} = EM\ket{\psi} = E\ket{\psi}$, which means that this error cannot be detected. So we define $\mathcal{N}(\mathcal{S}) = \{E \in \Pi^n \colon EM = ME, \ \forall M \in \mathcal{S}\}$ which can be seen as the set of errors which cannot be detected by the code $\mathcal{T}(\mathcal{S})$. We can also define the distance of a code, which is the smallest weight of an operator in the set $\mathcal{N}(\mathcal{S}) \setminus \mathcal{S}$, where the weight is the number of elements in the tensor product which compose the operator that are not the identity.\\
	 Finally the last question to ask is, how do we derive circuits from the stabilizer? The idea is to use circuits like the following:
	 \begin{figure}[H]
		\centering
		\begin{quantikz}
		\lstick{$\ket{0}$} & \gate{H} & \ctrl{1} & \gate{H} & \meter{} & \cw \\ 
			\lstick{$\ket{\psi}$} & \qw & \gate{U} & \qw & \qw & \qw
		\end{quantikz}
	 \end{figure}
	 \[ \ket{0}\ket{\psi} \xrightarrow{H\otimes I} \frac{1}{\sqrt{2}}\Big[\ket{0}\ket{\psi}+ \ket{1}\ket{\psi}\Big] \xrightarrow{I\otimes U} \frac{1}{\sqrt{2}}\Big[\ket{0}\ket{\psi}+ \ket{1}U\ket{\psi}\Big] \xrightarrow{H\otimes I} \frac{1}{2}\Big[(\ket{0}+\ket{1})\ket{\psi}+ (\ket{0}-\ket{1})U\ket{\psi}\Big] \]
	 \[ = \frac{1}{2}\ket{0}(\ket{\psi} + U\ket{\psi}) + \frac{1}{2}\ket{1}(\ket{\psi} - U\ket{\psi}) \]
	 If $\ket{\psi}$ is a +1 eigenstate of $U$ then it will only remain the term $\ket{0}\ket{\psi}$ and we'll measure 0 on the fist qubit, while if is a -1 eigenstate of $U$ it will only remain the term $\ket{1}\ket{\psi}$ and we'll measure 1 on the first qubit. In any case we recover the input state $\ket{\psi}$ on the second register and we get of the first register the information of the sign of the eigenvalue associated to the eigenstate $\ket{\psi}$.\\
	 So if $U$ is a stabilizer of $\ket{\psi}$ we'll always measure 0. In instead $\psi$ has been affected by some error $E$ correctable from $\mathcal{T}(S)$, hence $UE\ket{\psi} = -EU\ket{\psi}  = -E\ket{\psi}$, we'll measure 1, and therefore we know that an error occurred.\\
	 So what we  can do is build a series of circuits like the following one for generator of the stabilizer code. We'll se an example of this with the 5-qubit encoding. 
	 
	 \subsection{Five qubit code}
	 The five qubit code is the smallest code (in term of number of qubits for the encoding) which is able to detect bit-phase flip errors on a single qubit, using the following encoding:
	 \[ \ket{0} \rightarrow \ket{00000} + \ket{00000} + \ket{00000} + \ket{00000} + \ket{00000} + \ket{00000} + \ket{00000} + \ket{00000} + \]
	 \[ \ket{00000} + \ket{00000} + \ket{00000} + \ket{00000} + \ket{00000} + \ket{00000} + \ket{00000} + \ket{00000} \]
	 \[ \ket{1} \rightarrow \ket{00000} + \ket{00000} + \ket{00000} + \ket{00000} + \ket{00000} + \ket{00000} + \ket{00000} + \ket{00000} + \]
	 \[ \ket{00000} + \ket{00000} + \ket{00000} + \ket{00000} + \ket{00000} + \ket{00000} + \ket{00000} + \ket{00000} \]
	 The generators of the code are the following:
	 \[
	 \begin{array}{c|c c c c c}
	 	g_1 & X & Z & Z & X & I\\
	 	g_2 & I & X & Z & Z & X\\
	 	g_3 & X & I & X & Z & Z\\
	 	g_4 & Z & X & I & X & Z\\
	 \end{array}
	 \]
	 Once we have figured out a circuit for the encoding (which is not particularly interesting) what we'll do is applying a series of measurements:
	 \begin{figure}[H]
	 	\centering
	 	\resizebox{10cm}{!}{%
	 	\begin{quantikz}
	 		\lstick{$\ket{\psi}$} & \qwbundle{5} & \gate{XZZXI} & \gate{IXZZX} & \gate{XIXZZ} & \gate{ZXIXZ} & \qw & \qw \\
	 		\lstick{$\ket{0}$} & \gate{H} & \ctrl{-1} & \qw & \qw & \qw & \gate{H} & \meter{} & \cw \\ 
	 		\lstick{$\ket{0}$} & \gate{H} & \qw & \ctrl{-2} & \qw & \qw & \gate{H} & \meter{} & \cw \\
	 		\lstick{$\ket{0}$} & \gate{H} & \qw & \qw & \ctrl{-3} & \qw & \gate{H} & \meter{} & \cw \\ 
	 		\lstick{$\ket{0}$} & \gate{H} & \qw & \qw & \qw & \ctrl{-4} & \gate{H} & \meter{} & \cw \\
	 	\end{quantikz}
 		}
	 \end{figure}
 	note that $XIXZZ$ is a shorthand notation for $X\otimes I\otimes X\otimes Z\otimes Z$
	 
	 
	 
	 
	\newpage
	\bibliographystyle{unsrt}
	\bibliography{bibliography}

\end{document} 
